Ένα υπολογιστικό μοντέλο μπορεί να θεωρηθεί αξιόπιστο μόνο μετά από επικύρωση με πειραματικά δεδομένα. Επομένως, έχει ιδιαίτερο νόημα η δυνατότητα προσαρμογής του μοντέλου σε πειραματικές μετρήσεις αλλά και ο έλεγχος ικανότητας του μοντέλου να γενικεύει σε σύνολο δεδομένων πέρα από τα δεδομένα βαθμονόμησης. Παρακάτω παρουσιάζονται συνοπτικά οι συνήθεις τεχνικές μέτρησης της θερμοκρασίας και της φθοράς του πέλματος, καθώς και η διαδικασία προσδιορισμού των παραμέτρων του μοντέλου.

\subsection{Πειραματικές μετρήσεις}

Η θερμοκρασία του πέλματος καταγράφεται σε πραγματικό χρόνο κατά τη διάρκεια του αγώνα, μέσω δύο κυρίως συστημάτων αισθητήρων. Το πρώτο αποτελείται από αισθητήρες υπερύθρων (\textit{infrared sensors}) τοποθετημένους σταθερά στο πλαίσιο του οχήματος με στόχευση στην επιφάνεια του ελαστικού, οι οποίοι μετρούν την κατανομή της επιφανειακής θερμοκρασίας κατά το πλάτος του πέλματος. Το δεύτερο είναι ένα εσωτερικό σύστημα παρακολούθησης πίεσης και θερμοκρασίας ελαστικού (\textit{TPMS}), τοποθετημένο εντός της ζάντας, από όπου λαμβάνεται η θερμοκρασία και πίεση του εγκλωβισμένου αέρα. Στην παρούσα εργασία, ως μέγεθος σύγκρισης με το μοντέλο θεωρείται η μέση θερμοκρασία της επιφάνειας του πέλματος.\\[10pt]

Η φθορά, αντίθετα, δε μετράται συνεχώς κατά τη διάρκεια του αγώνα. Ο συνήθης τρόπος προσδιορισμού της είναι η μέτρηση του πάχους του πέλματος πριν και μετά από μια περίοδο συνεχόμενης χρήσης (\textit{stint}), σε πολλαπλά σημεία της περιφέρειας και του πλάτους του ελαστικού. Ως καμπύλη αναφοράς λαμβάνεται επομένως η συνολική συσσωρευμένη φθορά ανά τροχό στο τέλος της περιόδου \cite{suttill2025treadwear}.

\subsection{Προσδιορισμός παραμέτρων}

Οι παράμετροι που υπόκεινται σε βαθμονόμηση είναι οι $p_{1\dots6}$ του θερμοδυναμικού μοντέλου (πίνακας \ref{tab:params_thermal}) και οι παράμετροι του μοντέλου φθοράς (πίνακας \ref{tab:params_wear}). Ο προσδιορισμός γίνεται ξεχωριστά ανά τροχό, ώστε να αποτυπωθούν ασυμμετρίες στις παραμέτρους. Ωστόσο, θεωρητικά οι περισσότερες παράμετροι, και ιδιαίτερα οι παράμετροι φθοράς, είναι ίσες μεταξύ των ελαστικών, με εξαίρεση τις παραμέτρους συναγωγής οι οποίες μεταβάλλονται λόγω της διαφοράς στη ροή γύρω από κάθε τροχό. \\[10pt]

Ως αντικειμενική συνάρτηση επιλέγεται το τετραγωνικό μέσο σφάλμα (\textit{RMSE}) μεταξύ της εξόδου του μοντέλου και της καμπύλης αναφοράς:

\begin{equation}
    \mathrm{RMSE}(\mathbf{p}) = \sqrt{\frac{1}{N}\sum_{i=1}^{N}\bigl(y_{\text{μοντ.},i}(\mathbf{p}) - y_{\text{αναφ.},i}\bigr)^2}
    \label{eq:paramFittingObjective}
\end{equation}

όπου $\mathbf{p}$ το διάνυσμα των υπό εκτίμηση παραμέτρων, $y_{\text{μοντ.}}(\mathbf{p})$ και $y_{\text{αναφ.}}$ οι τιμές του μοντέλου και της αναφοράς αντίστοιχα (θερμοκρασία ή φθορά), και $N$ ο αριθμός σημείων δειγματοληψίας. Η ελαχιστοποίηση πραγματοποιείται με τη συνάρτηση \textit{lsqnonlin} της \textit{Matlab} \cite{mathworks_lsqnonlin}, με αρχική εκτίμηση τις τιμές των πινάκων (\ref{tab:params_thermal}) και (\ref{tab:params_wear}).
